\section{Постановка задачи}

Требуется разработать библиотеку \texttt{Compio}, обеспечивающую прозрачное сжатие данных с возможностью произвольного доступа на чтение, запись, вставку и удаление. Библиотека должна предоставлять интерфейс, аналогичный стандартным потокам ввода-вывода \texttt{stdio}, и поддерживать хранение нескольких логических файлов в одном физическом файле-контейнере.

\subsection{Функциональные требования}

\begin{enumerate}
    \item \textbf{Интерфейс}. Библиотека должна предоставлять функции для работы с архивом и файлами внутри него:
    \begin{itemize}
        \item открытие и закрытие архива;
        \item открытие и закрытие логического файла по имени;
        \item чтение данных из текущей позиции;
        \item запись данных в текущую позицию (с возможным перекрытием существующих данных);
        \item вставка данных в произвольную позицию (с увеличением размера файла);
        \item удаление данных из произвольной позиции (с уменьшением размера файла);
        \item позиционирование (аналог \texttt{seek}) и получение текущей позиции (аналог \texttt{tell}).
    \end{itemize}
    Имена функций и константы должны иметь единый префикс (например, \texttt{compio\_}), а семантика операций должна быть максимально близка к \texttt{stdio}, чтобы упростить миграцию существующего кода.

    \item \textbf{Прозрачное сжатие}. Сжатие должно выполняться автоматически при записи данных, а разжатие — при чтении. Пользователь не должен явно вызывать операции сжатия или разжатия.

    \item \textbf{Произвольный доступ}. Библиотека должна обеспечивать возможность чтения, записи, вставки и удаления данных в любой позиции логического файла без необходимости полной декомпрессии всего архива.

    \item \textbf{Многофайловость}. Один физический файл-контейнер должен вмещать произвольное количество логических файлов. Имена файлов должны быть строковыми идентификаторами (ограниченной длины). Должна быть возможность перечислить имена файлов, получить размер каждого файла.

    \item \textbf{Модульность алгоритма сжатия}. Библиотека должна поддерживать выбор алгоритма сжатия (Zlib, LZ4, Zstd, Brotli) как при создании нового архива, так и при открытии существующего. Должна быть предусмотрена возможность подключения пользовательского алгоритма сжатия через предоставление трёх функций: \texttt{compress}, \texttt{decompress}, \texttt{get\_buf\_size}. Тип использованного алгоритма должен сохраняться в контейнере, чтобы при открытии архива библиотека могла проверить согласованность.

    \item \textbf{API}. Библиотека должна предоставлять C-совместимый интерфейс (стабильный ABI) для использования из языков C и C++, а также из других языков через механизмы FFI. При желании может быть предоставлена дополнительная удобная C++ обёртка.

    \item \textbf{Кроссплатформенность}. Библиотека должна компилироваться и работать под операционными системами Linux и Windows.

    \item \textbf{Самодостаточность}. Весь архив должен храниться в одном файле файловой системы. Для чтения и модификации архива не требуется никаких дополнительных файлов. Файл должен быть переносимым между разными машинами (при условии одинаковой архитектуры и используемых кодеков).
\end{enumerate}

\subsection{Нефункциональные требования}

\begin{enumerate}
    \item \textbf{Эффективность при локальном доступе}. При высокой локальности обращений (повторные чтения или записи в одних и тех же областях) пропускная способность должна быть сопоставима или превосходить \texttt{stdio} благодаря кэшированию. При случайном доступе допустимо некоторое снижение производительности из-за накладных расходов на индексацию и сжатие.

    \item \textbf{Эффективность операций вставки и удаления}. Вставка и удаление должны выполняться за время, логарифмически зависящее от количества блоков в файле (амортизированно), без необходимости перепаковки всего файла.

    \item \textbf{Управление фрагментацией}. Библиотека должна контролировать внутреннюю фрагментацию (появление слишком маленьких или слишком больших блоков) и внешнюю фрагментацию (пустоты в файле-контейнере). Должны быть предусмотрены механизмы для поддержания приемлемой степени фрагментации (например, через выбор стратегии аллокатора или периодическую дефрагментацию).

    \item \textbf{Кэширование}. Библиотека должна использовать кэш узлов индекса и кэш распакованных данных с ограничением объёма используемой оперативной памяти. Алгоритм вытеснения должен быть предсказуемым (например, LRU).

    \item \textbf{Отсутствие внешних зависимостей}. Библиотека не должна требовать установки дополнительных библиотек во время выполнения (кроме стандартной библиотеки языка C++ и, опционально, выбранных кодеков, которые могут быть статически слинкованы или предоставлены пользователем).

    \item \textbf{Документация}. API должен быть документирован (например, с помощью Doxygen), а репозиторий с исходным кодом — открытым.

    \item \textbf{Тестирование}. Библиотека должна быть покрыта модульными тестами, проверяющими как отдельные компоненты (B-дерево, кэш, аллокатор), так и сквозные сценарии использования (включая случайные последовательности операций).
\end{enumerate}

\subsection{Границы применимости}

Библиотека не предназначена для многопоточного доступа к одному архиву без внешней синхронизации (каждая операция предполагается выполняемой в одном потоке, либо пользователь должен обеспечивать блокировки). Не гарантируется работа с файлами, размер которых превышает ёмкость 64-битных адресов (более $2^{64}$ байт), однако для практических задач это ограничение несущественно. Библиотека не реализует собственное шифрование или проверку целостности на уровне отдельных операций (кроме контрольной суммы блока), полная проверка целостности архива не поддерживается.